\documentclass[a4paper, 10pt]{jarticle}
\usepackage[dvipdfmx]{graphicx}
\usepackage{url}
\usepackage{amsmath}
\usepackage{latexsym}
\usepackage{tikz}
\usepackage{float}
\usepackage{subcaption}
\usepackage{amssymb}

\setlength{\textwidth}{165mm}
\setlength{\marginparwidth}{40mm}
\setlength{\textheight}{225mm}
\setlength{\topmargin}{-5mm}
\setlength{\oddsidemargin}{-3.5mm}

\begin{document}
\title{教科書確認用}
\maketitle

% 練習1
下の図において、$\sin\theta$、$\cos\theta$、$\tan\theta$の値を、それぞれ求めよ。

% 練習2
次の値を求めよ。
\begin{enumerate}
  \item $\cos30^\circ,\ \tan30^\circ$
  \item $\sin45^\circ,\ \tan45^\circ$
  \item $\sin60^\circ,\ \cos60^\circ$
\end{enumerate}

% 練習3
次の値を、三角比の表を用いて求めよ。
\begin{enumerate}
  \item $\sin22^\circ$
  \item $\cos68^\circ$
  \item $\tan86^\circ$
\end{enumerate}

% 練習4
次のような鋭角$\theta$のおよその大きさを、三角比の表を用いて求めよ。
\begin{enumerate}
  \item $\sin\theta=0.26$
  \item $\cos\theta=0.53$
  \item $\tan\theta=4.3$
\end{enumerate}

% 練習5
傾斜角$20^\circ$の坂をまっすぐに$500\,\mathrm{m}$登るとき、
鉛直方向に何$\mathrm{m}$上がることになるか。
また、水平方向に何$\mathrm{m}$進むことになるか。
$1\,\mathrm{m}$未満は四捨五入して求めよ。

% 練習6
目の高さが$1.6\,\mathrm{m}$の地点から塔の先端を見上げたところ、
仰角は$40^\circ$であった。
塔の高さを求めよ。ただし、小数第2位を四捨五入せよ。

% 練習7
右の道路標識は道路の勾配を表したもので、
水平距離に対する高さの割合が$10\%$、
すなわち水平距離$100\,\mathrm{m}$に対して
$10\,\mathrm{m}$の割合で高くなることを示している。
この道路の傾斜角を求めよ。

% 練習8
$\theta$は鋭角とする。
$\cos\theta=\dfrac{2}{3}$のとき、
$\sin\theta$と$\tan\theta$の値を求めよ。

% 練習9
$\theta$は鋭角とする。
$\tan\theta=\dfrac{1}{2}$のとき、
$\sin\theta$と$\cos\theta$の値を求めよ。

% 練習10
次の三角比を求めよ。
\begin{enumerate}
  \item $\sin120^\circ$
  \item $\cos135^\circ$
  \item $\tan150^\circ$
\end{enumerate}

% 練習11
図を用いて、次の等式が成り立つことを確かめよ。
\begin{enumerate}
  \item $\cos\left(90^\circ-\theta\right)=\sin\theta$
  \item $\tan\left(90^\circ-\theta\right)=\dfrac{1}{\tan\theta}$
\end{enumerate}

% 練習12
下の図において、次の三角比を求めよ。
\begin{enumerate}
  \item $\sin\theta$
  \item $\cos\theta$
  \item $\tan\theta$
\end{enumerate}

% 練習13
次の値を求めよ。
\begin{enumerate}
  \item $\sin120^\circ$
  \item $\cos135^\circ$
  \item $\tan120^\circ$
\end{enumerate}

% 練習14
$\theta$は鈍角とする。
次の等式を満たす$\theta$を求めよ。
\begin{enumerate}
  \item $\sin\theta=\dfrac{1}{2}$
  \item $\cos\theta=-\dfrac{\sqrt{3}}{2}$
  \item $\cos\theta=-\dfrac{1}{2}$
  \item $\sin\theta=\dfrac{\sqrt{3}}{2}$
\end{enumerate}

% 練習15
$\theta$は鈍角とする。
$\tan\theta=-\dfrac{1}{2}$のとき、
$\sin\theta$と$\cos\theta$の値を求めよ。

% 練習16
$0^\circ\leqq\theta\leqq180^\circ$とする。
$\sin\theta$、$\cos\theta$、$\tan\theta$のうち
1つが次の値をとるとき、他の2つの値を求めよ。
\begin{enumerate}
  \item $\sin\theta=\dfrac{2}{3}$
  \item $\cos\theta=\dfrac{1}{4}$
\end{enumerate}

% 練習17
$0^\circ\leqq\theta\leqq180^\circ$とする。
$\tan\theta$が次の値をとるとき、
$\sin\theta$と$\cos\theta$の値を求めよ。
\begin{enumerate}
  \item $\tan\theta=-\sqrt{2}$
  \item $\tan\theta=\dfrac{1}{3}$
\end{enumerate}

% 練習18
次の直線と$x$軸の正の向きとのなす角を求めよ。
\begin{enumerate}
  \item $y=x$
  \item $y=-\sqrt{3}x$
\end{enumerate}

% 練習19
$\triangle ABC$において、次のものを求めよ。
\begin{enumerate}
  \item $a=6,\ B=30^\circ,\ C=75^\circ$のとき$b$
  \item $c=4,\ A=30^\circ,\ B=45^\circ$のとき$a$
\end{enumerate}

% 練習20
$\triangle ABC$において、外接円の半径を$R$とする。
次のものを求めよ。
\begin{enumerate}
  \item $b=8,\ B=60^\circ$のとき$R$
  \item $a=R$のとき$A$
\end{enumerate}

% 練習21
$\triangle ABC$において、次のものを求めよ。
\begin{enumerate}
  \item $a=3,\ c=2\sqrt{2},\ B=45^\circ$のとき$b$
  \item $a=8,\ b=7,\ C=120^\circ$のとき$c$
\end{enumerate}

% 練習22
$\triangle ABC$において、
$a=\sqrt{5},\ b=\sqrt{2},\ A=45^\circ$のとき、
$c$を求めよ。

% 練習23
$\triangle ABC$において、
$a=2,\ b=7,\ c=3\sqrt{3}$のとき、
$B$を求めよ。

% 練習24
$\triangle ABC$において、
$a=3,\ b=5,\ c=6$のとき、
$A,\ B,\ C$は鋭角、直角、鈍角のいずれであるか。

% 練習25
$\triangle ABC$において、3辺の長さが次のようなとき、
$\triangle ABC$は鋭角三角形、直角三角形、鈍角三角形の
いずれであるか。
\begin{enumerate}
  \item $a=4,\ b=3,\ c=2$
  \item $a=5,\ b=6,\ c=7$
\end{enumerate}

% 練習26
$\triangle ABC$において、
$b=2,\ c=\sqrt{6}+\sqrt{2},\ A=45^\circ$のとき、
$a,\ B,\ C$を求めよ。

% 練習27
$\triangle ABC$において、
$b=2,\ c=\sqrt{2},\ C=30^\circ$のとき、
$a,\ A,\ B$を求めよ。

% 練習28
$\triangle ABC$において、次の等式が成り立つとき、
この三角形の最も小さい角の余弦の値を求めよ。
\[
\frac{\sin A}{8}=\frac{\sin B}{7}=\frac{\sin C}{3}
\]

% 練習29
次のような$\triangle ABC$の面積$S$を求めよ。
\begin{enumerate}
  \item $a=3,\ c=8,\ B=60^\circ$
  \item $a=\sqrt{6},\ b=2\sqrt{2},\ C=150^\circ$
\end{enumerate}

% 練習30
$\triangle ABC$において、
$a=4,\ b=3,\ c=2$のとき、
この三角形の面積$S$を求めよ。

% 練習31
$\triangle ABC$において、
$\mathrm{AB}=5,\ \mathrm{AC}=3,\ A=120^\circ$とし、
$\angle A$の二等分線と辺$\mathrm{BC}$の交点を$D$とする。
線分$\mathrm{AD}$の長さを求めよ。

% 練習32
円に内接する四角形$\mathrm{ABCD}$において、
$\mathrm{AB}=5,\ \mathrm{BC}=8,\ \mathrm{CD}=3,\ \angle B=60^\circ$
とするとき、四角形$\mathrm{ABCD}$の面積$S$を求めよ。

% 練習33
$a=7,\ b=5,\ c=4$である$\triangle ABC$の
内接円の半径$r$を求めよ。

% 練習34
右の図のような、
$\mathrm{AB}=6,\ \mathrm{AD}=4,\ \mathrm{AE}=3$
である直方体$\mathrm{ABCD-EFGH}$がある。
$\triangle DEG$の面積$S$を求めよ。

% 練習35
1辺の長さが$a$の正四面体$\mathrm{ABCD}$において、
辺$\mathrm{CD}$の中点を$M$とする。
このとき、次のものを求めよ。
\begin{enumerate}
  \item $\cos\angle ABM$の値
  \item $\triangle ABM$の面積
\end{enumerate}

% 練習36
$\mathrm{PA}=\mathrm{PB}=\mathrm{PC}=4$、
$\mathrm{AB}=6,\ \mathrm{BC}=4,\ \mathrm{CA}=5$
である三角錐$\mathrm{PABC}$の体積$V$を求めよ。

% 練習37
テレビ塔が立っている地点$H$と同じ標高の地点$A$から、
テレビ塔の先端$P$を見たところ、仰角が$60^\circ$であった。
また、$A$から$40\,\mathrm{m}$離れた地点$B$では、
$\angle HAB=15^\circ,\ \angle HBA=135^\circ$であった。
テレビ塔の高さ$\mathrm{PH}$を求めよ。

\end{document}

